\section{Анализ подходов к верификации подсистемы памяти}
\begin{frame}[plain, noframenumbering]
  \begin{center}
    \Huge
    Анализ подходов~~к~~верификации~~подсистемы~~памяти
  \end{center}
\end{frame}

\begin{frame}
  \frametitle{Актуальность}
  \begin{itemize}
    \item Рост числа ядер в~многопроцессорных системах существенно
      усложняет архитектуру подсистемы памяти;
    \item Усложнение иерархии памяти и~увеличение параллелизма требуют
      развития механизмов синхронизации;
    \item Верификация консистентности памяти --- корректности результатов
      операций чтения/записи при параллельном исполнении ---
      остаётся одной из~наиболее сложных задач;
    \item Нарушения спецификации модели памяти приводят к~трудно
      воспроизводимым сбоям.
  \end{itemize}
\end{frame}

\begin{frame}
  \frametitle{Модели памяти}
  Модель памяти --- часть архитектуры процессора, определяющая порядок,
  в~котором результаты операций из разных потоков становятся видимыми
  друг для друга:
  \begin{description}
    \item[SC] --- последовательная консистентность (strongest);
    \item[TSO] --- последовательная консистентность с~буферизацией
      записей;
    \item[XC] --- расширенная консистентность для слабоупорядоченных
      систем;
    \item[RC] --- слабая консистентность (release/acquire);
    \item[RVWMO] --- слабая модель памяти RISC-V.
  \end{description}
\end{frame}

\begin{frame}
  \frametitle{Обзор существующих подходов}
  \begin{itemize}
    \item \textbf{Литмус-тестирование} --- ориентировано на~проверку
      соответствия системы заданной модели памяти, тесты содержат
      небольшое число операций;
    \item \textbf{Стресс-тестирование} --- тысячи конкурирующих
      обращений, но~экспоненциальный рост числа трасс;
    \item \textbf{Pre-silicon верификация} --- направленное тестирование
      на~ранних стадиях, но~недостаточная масштабируемость;
    \item \textbf{Аппаратные механизмы трассировки} --- требуют
      дополнительных ресурсов на~кристалле.
  \end{itemize}
\end{frame}

\begin{frame}
  \frametitle{Постановка задачи}
  Свести задачу анализа консистентности памяти к~построению всех
  линейных порядков частично упорядоченного множества операций
  с~памятью, учитывая модель памяти и~классы эквивалентности
  перестановок.
  \medskip

  Для достижения цели требовалось:
  \begin{enumerate}
    \item свести многопоточную программу к~частично упорядоченному
      множеству операций;
    \item разработать алгоритмы построения линейных порядков с~учётом
      эквивалентных перестановок;
    \item разработать граф состояний для компактного описания результатов;
    \item разработать алгоритмы верификации соответствия аппаратуры
      модели памяти.
  \end{enumerate}
\end{frame}

\section{Представление программы в~виде ЧУМ}
\begin{frame}[plain, noframenumbering]
  \begin{center}
    \Huge
    Представление~~программы~~в~~виде~~ЧУМ
  \end{center}
\end{frame}

\begin{frame}
  \frametitle{Математическая модель}
  \begin{block}{Представление программы}
    Многопоточная программа представляется в~виде частично упорядоченного
    множества (ЧУМ, partial~order) операций с~памятью, построенного
    с~учётом выбранной модели памяти.
  \end{block}
  \begin{itemize}
    \item Каждая операция с~памятью (загрузка, сохранение) --- элемент
      множества;
    \item Порядок задаётся программным порядком и~ограничениями модели
      памяти;
    \item Модель памяти задаёт множество допустимых результатов
      исполнения.
  \end{itemize}
\end{frame}

\begin{frame}
  \frametitle{Учёт классов эквивалентности}
  \begin{itemize}
    \item Операции, не~связанные зависимостями, могут быть
      переупорядочены --- порождают классы эквивалентности перестановок;
    \item Учёт независимости и~зависимости операций сокращает
      пространство анализируемых линейных порядков;
    \item Снимаются ограничения на~типы операций и~пересечение адресов,
      присущие литмус-тестам;
    \item Модель накладывает лишь ограничения, задаваемые моделью
      памяти.
  \end{itemize}
\end{frame}

\section{Линейные порядки ЧУМ}
\begin{frame}[plain, noframenumbering]
  \begin{center}
    \Huge
    Линейные~~порядки~~ЧУМ
  \end{center}
\end{frame}

\begin{frame}
  \frametitle{Комбинаторный взрыв}
  Прямое построение всех линейных порядков ЧУМ сталкивается
  с~комбинаторным взрывом:
  \begin{itemize}
    \item в~худшем случае число порядков растёт как $n!$;
    \item средняя асимптотическая сложность прямого алгоритма ---
      $O(2^n)$.
  \end{itemize}
  Это обосновывает необходимость разработки специализированных
  алгоритмов и~структур данных.
\end{frame}

\begin{frame}
  \frametitle{Алгоритмы топологической сортировки}
  Разработаны алгоритмы построения линейных порядков:
  \begin{itemize}
    \item топологическая сортировка с~учётом \textbf{независимости}
      операций;
    \item топологическая сортировка с~учётом \textbf{зависимостей};
    \item использование заранее построенных таблиц перестановок.
  \end{itemize}
  \begin{alertblock}{Оценка сложности}
    В~предельном случае высокой независимости операций вычислительные
    затраты снижаются до~$O(n^2)$.
  \end{alertblock}
\end{frame}

\section{Компактное представление результатов}
\begin{frame}[plain, noframenumbering]
  \begin{center}
    \Huge
    Компактное~~представление~~результатов
  \end{center}
\end{frame}

\begin{frame}
  \frametitle{Граф состояний (State Graph)}
  \begin{block}{Структура данных}
    Граф состояний --- компактное и~полное представление всех возможных
    результатов исполнения многопоточной программы, позволяющее
    восстанавливать линейные порядки без их~явного хранения.
  \end{block}
  \begin{itemize}
    \item Позволяет преодолеть экспоненциальный рост числа состояний;
    \item Даёт возможность верификации программ с~большим числом
      конкурирующих обращений;
    \item Применим для программ с~десятками тысяч операций.
  \end{itemize}
\end{frame}

\begin{frame}
  \frametitle{Метод кластеризации}
  \begin{itemize}
    \item Предложен метод \textbf{кластеризации графа состояний}
      для декомпозиции задачи верификации на~независимые подзадачи;
    \item Введена \textbf{метрика насыщения}, позволяющая различать
      эффекты, обусловленные переупорядочиванием зависимых операций
      и~моделью памяти;
    \item Метрика насыщения позволяет останавливать исполнение программы
      при достижении полного покрытия кластеров состояний.
  \end{itemize}
\end{frame}

\begin{frame}
  \frametitle{Алгоритмы верификации}
  Предложены два подхода к~верификации соответствия аппаратуры
  заявленной модели памяти:
  \begin{enumerate}
    \item на основе построения всех линейных порядков частично
      упорядоченного множества;
    \item на основе анализа графа состояний.
  \end{enumerate}
  \medskip
  Доказана эквивалентность данных подходов к~верификации.
\end{frame}

\section{Результаты}
\begin{frame}[plain, noframenumbering]
  \begin{center}
    \Huge
    Результаты
  \end{center}
\end{frame}

\begin{frame}
  \frametitle{Экспериментальная апробация}
  Апробация проводилась на~тестовых сценариях, моделирующих:
  \begin{itemize}
    \item классические \textbf{литмус-тесты};
    \item нагрузки с~большим числом конкурирующих обращений
      к~памяти.
  \end{itemize}
  \begin{alertblock}{Масштабируемость}
    Предложенный подход позволяет масштабировать верификацию на~программы
    с~десятками тысяч операций, что выходит за~рамки применимости
    классических подходов.
  \end{alertblock}
\end{frame}

\begin{frame}
  \frametitle{Основные результаты работы}
  \begin{enumerate}
    \item Предложена модель представления программы в~виде ЧУМ операций
      с~памятью с~учётом моделей памяти (SC, TSO, XC, RC, RVWMO);
    \item Разработаны алгоритмы построения линейных порядков с~учётом
      классов эквивалентности перестановок;
    \item Разработана структура данных ``граф состояний'' для
      компактного и~полного представления результатов;
    \item Предложены алгоритмы верификации на~основе линейных порядков
      и~на~основе анализа графа состояний;
    \item Созданы метод кластеризации графа состояний и~метрика
      насыщения.
  \end{enumerate}
\end{frame}

\begin{frame}
  \frametitle{Перспективы дальнейшей работы}
  \begin{itemize}
    \item Адаптация подхода к~гетерогенным вычислительным архитектурам
      (GPU, FPGA, тензорные блоки) --- построение композиционной модели
      памяти через точки синхронизации;
    \item Развитие эвристик анализа на~основе метрики насыщения;
    \item Разработка адаптивных стратегий генерации тестовых сценариев,
      целенаправленно увеличивающих покрытие редких состояний.
  \end{itemize}
\end{frame}
